% 第2章 2.4.1：ℓ^2 标准单位向量 e_n 弱收敛到零而不范数收敛。
% (a) 固定测试向量 y_k=1/k（蓝色柱）与 e_6（黄色柱）；读数 <y,e_6>=y_6=1/6。
% (b) 读数 <y,e_n>=1/n 随 n 趋于零，而 ‖e_n‖=1 不变。
\begin{tikzpicture}[font=\figurefont\footnotesize]
\pgfplotsset{every axis/.style={width=56mm,height=36mm,axis lines=left,
  x axis line style={draw=FigureStroke,line width=.6pt},y axis line style={draw=FigureStroke,line width=.6pt},
  tick style={draw=FigureStroke},tick label style={font=\figurefont\scriptsize,text=FigureInk},
  label style={font=\figurefont\scriptsize},
  xlabel style={at={(axis description cs:1,0)},anchor=north west,yshift=1mm},
  ymajorgrids,grid style={draw=FigureGrid,line width=.4pt},
  title style={font=\figurefont\footnotesize},
  legend style={draw=none,fill=none,font=\figurefont\scriptsize,cells={anchor=west}}}}
\begin{axis}[name=A,ybar,bar width=1.6mm,xlabel={$k$},xmin=0.3,xmax=10.7,ymin=0,ymax=1.15,
  xtick={1,2,...,10},ytick={0,0.5,1},title={(a) 测试向量 $y$ 与 $e_6$},
  legend style={at={(0.98,0.8)},anchor=north east}]
\addplot[fill=FigureBlueFill,draw=FigureBlue,bar shift=-0.9mm] coordinates {(1,1) (2,0.5) (3,0.3333) (4,0.25) (5,0.2) (6,0.1667) (7,0.1429) (8,0.125) (9,0.1111) (10,0.1)};
\addplot[fill=FigureYellow,draw=FigureYellow,bar shift=0.9mm] coordinates {(6,1)};
\legend{$y_k=1/k$,$e_6$}
\end{axis}
\begin{axis}[at={(A.east)},anchor=west,xshift=10mm,xlabel={$n$},xmin=0,xmax=20.5,ymin=0,ymax=1.15,
  xtick={1,5,10,15,20},ytick={0,0.5,1},title={(b) 读数与范数},
  legend style={at={(0.98,0.55)},anchor=east}]
\addplot[draw=FigureBlue,line width=1.1pt,mark=*,mark size=1pt] coordinates {(1,1) (2,0.5) (3,0.3333) (4,0.25) (5,0.2) (6,0.1667) (7,0.1429) (8,0.125) (9,0.1111) (10,0.1) (11,0.0909) (12,0.0833) (13,0.0769) (14,0.0714) (15,0.0667) (16,0.0625) (17,0.0588) (18,0.0556) (19,0.0526) (20,0.05)};
\addplot[draw=FigureYellow,line width=1.1pt,dash dot] coordinates {(1,1) (20,1)};
\legend{{$\langle y,e_n\rangle$},{$\|e_n\|$}}
\end{axis}
\end{tikzpicture}
